\section{Frontend (5123101)}
Im Rahmen der Frontend-Entwicklung wurden verschiedene moderne Webtechnologien eingesetzt. Die Auswahl der verwendeten Technologien erfolgte bewusst unter Berücksichtigung von Wartbarkeit, Performance, Skalierbarkeit sowie Entwicklerfreundlichkeit.

\subsection{Auswahl des Frameworks}
Für die Entwicklung der Benutzeroberfläche wurde React verwendet. Gegenüber alternativen Frameworks wie Vue.js oder Angular bietet React mehrere Vorteile.

Ein wesentlicher Vorteil besteht in der komponentenbasierten Architektur. Die Benutzeroberfläche kann dadurch in kleine, wiederverwendbare Komponenten aufgeteilt werden. Dies verbessert sowohl die Wartbarkeit als auch die Strukturierung des Projekts erheblich.

Darüber hinaus zeichnet sich React durch eine hohe Flexibilität aus. Im Gegensatz zu Angular macht React nur wenige Vorgaben bezüglich Projektstruktur oder zusätzlicher Technologien. Dadurch können passende Bibliotheken und Werkzeuge gezielt ausgewählt und kombiniert werden.

Ein weiterer wichtiger Aspekt ist das große Ökosystem sowie die starke Verbreitung im professionellen Umfeld. Für React existieren zahlreiche Bibliotheken, Erweiterungen und Dokumentationen, wodurch moderne Entwicklungsstandards unterstützt werden und viele Probleme bereits gelöst sind.

Zusätzlich besitzt React eine vergleichsweise geringe Einstiegshürde und bietet den Vorteil, dass im Team bereits Vorkenntnisse aus einem anderen Modul vorliegen.

\subsection{Verwendete Programmiersprache}
Das Frontend wurde bewusst mit TypeScript anstelle von reinem JavaScript umgesetzt.

TypeScript erweitert JavaScript um statische Typisierung. Dadurch können viele Fehler bereits während der Entwicklung erkannt werden, bevor die Anwendung ausgeführt wird. Dies erhöht die Codequalität und reduziert typische Laufzeitfehler.

Gerade bei größeren Anwendungen verbessert TypeScript außerdem die Wartbarkeit des Codes. Funktionen, Komponenten und Datenstrukturen besitzen klar definierte Typen, wodurch Schnittstellen innerhalb der Anwendung eindeutig beschrieben werden können.

Ein weiterer Vorteil liegt in der besseren Unterstützung durch moderne Entwicklungsumgebungen. Funktionen wie automatische Vervollständigung, Typenhinweise oder sichere Refactorings erleichtern die Entwicklung erheblich.

Da das Frontend mit unterschiedlichen Datenobjekten und API-Antworten arbeitet, bietet TypeScript zusätzliche Sicherheit bei der Verarbeitung dieser Daten.

\subsection{Styling mit Tailwind CSS}
Für das Design der Benutzeroberfläche wurde Tailwind CSS verwendet.

Tailwind verfolgt einen sogenannten \enquote{Utility-First}-Ansatz. Dabei werden fertige CSS-Klassen direkt innerhalb der Komponenten genutzt, anstatt umfangreiche eigene CSS-Dateien zu erstellen.

Ein großer Vorteil von Tailwind besteht in der schnellen Entwicklung moderner Benutzeroberflächen. Viele häufig benötigte Styles wie Abstände, Farben, Größen oder Flexbox-Layouts stehen bereits als Klassen zur Verfügung. Dadurch reduziert sich der Aufwand für eigenes CSS deutlich.

Zusätzlich sorgt Tailwind für eine konsistente Gestaltung der gesamten Anwendung. Durch wiederverwendbare Utility-Klassen entsteht ein einheitliches Designsystem, wodurch Inkonsistenzen vermieden werden.

Ein weiterer Vorteil ist die gute Integration in React-Komponenten. Styling und Komponentenlogik befinden sich nah beieinander, wodurch die Übersichtlichkeit verbessert wird.

\subsection{Build-Tool Vite}
Für die Entwicklung und den Build-Prozess wurde Vite eingesetzt.

Vite bietet im Vergleich zu klassischen Build-Werkzeugen wie Webpack deutlich schnellere Start- und Build-Zeiten. Besonders während der Entwicklung profitieren Entwickler von kurzen Ladezeiten und schnellem Hot Module Replacement (HMR). Änderungen am Code werden dabei nahezu sofort im Browser sichtbar.

Zusätzlich zeichnet sich Vite durch eine einfache Konfiguration aus. Viele Funktionen sind bereits standardmäßig integriert, wodurch weniger komplexe Konfigurationsdateien notwendig sind.

Darüber hinaus unterstützt Vite moderne Webstandards und ist optimal auf React sowie TypeScript abgestimmt.

\subsection{Routing mit React-Router}
Für die Navigation innerhalb der Anwendung wurde React-Router verwendet.

React-Router ermöglicht die Umsetzung einer Single-Page-Application (SPA), bei der Seitenwechsel ohne vollständiges Neuladen der Anwendung stattfinden. Dadurch entsteht eine flüssigere Benutzererfahrung.

Zusätzlich erlaubt React-Router eine klare Strukturierung der verschiedenen Seiten und Ansichten der Anwendung. Einzelne Bereiche können dadurch sauber voneinander getrennt werden.

Ein weiterer Vorteil ist die einfache Integration in React-Komponenten. Routen können flexibel definiert und mit Layouts oder geschützten Bereichen kombiniert werden.
